% Justificación
La epilepsia es una condición neurológica de alta prevalencia, cuyo diagnóstico depende en gran medida del análisis de señales electroencefalográficas (EEG). La interpretación manual de estos estudios multicanal genera una necesidad de inversión temporal por parte de los expertos debido a su alta experiencia y se ve limitada por la disponibilidad de recursos clínicos. Estos factores pueden retrasar la identificación oportuna de eventos relevantes. 

Aunque las técnicas de inteligencia artificial han mostrado avances significativos en la detección automática de eventos epilépticos, persisten limitaciones estructurales asociadas al entorno de evaluación predominante. Muchos modelos son entrenados y validados bajo condiciones controladas que no reflejan la variabilidad interpaciente ni las condiciones reales de práctica clínica. Como consecuencia, la generalización y estabilidad de los sistemas al enfrentarse a nuevos sujetos continúan siendo un desafío abierto.

La creciente complejidad de los modelos ha incrementado la opacidad en sus decisiones, dificultando su interpretación y capacidad de ser auditados. Este factor es esencial para el entorno clínico, donde la adopción de herramientas automatizadas requiere, no solo un alto desempeño, sino también transparencia y trazabilidad que permitan al especialista comprender y validar los resultados generados.

En este contexto, se vuelve necesario fortalecer el desarrollo de metodologías que prioricen la robustez, la generalización entre individuos y la posibilidad de auditoría clínica. El beneficio potencial de abordar este problema radica en contribuir a herramientas de apoyo al diagnóstico que sean técnicamente sólidas, clínicamente confiables y adecuadas para su implementación en entornos médicos reales.